\noindent\hspace*{2mm}%
\begin{minipage}[t]{\dimexpr\linewidth/2-4mm\relax}
\vspace{0pt}
\textbf{Versione quadratica}
\begin{unipdcode}[language=Python,
  basicstyle=\ttfamily\scriptsize,
  numbers=none,
  xleftmargin=0pt,
  xrightmargin=0pt,
  framexleftmargin=0.3em]
def differenza_massima_quadratica(valori):
    if not valori:                             # 1
        raise ValueError("Lista vuota")
    massimo = 0                                # 1
    for i in range(len(valori)):               # n iterazioni
        for j in range(len(valori)):           # n iterazioni
            differenza = valori[i] - valori[j] # 4 (2 letture, sottrazione, assegnamento)
            if differenza > massimo:           # 1
                massimo = differenza           # 1 se vero
    return massimo                             # escluso
\end{unipdcode}
\end{minipage}%
\hspace*{4mm}%
\begin{minipage}[t]{\dimexpr\linewidth/2-4mm\relax}
\vspace{0pt}
\textbf{Versione lineare}
\begin{unipdcode}[language=Python,
  basicstyle=\ttfamily\scriptsize,
  numbers=none,
  xleftmargin=0pt,
  xrightmargin=0pt,
  framexleftmargin=0.3em]
def differenza_massima_lineare(valori):
    if not valori:                      # 1
        raise ValueError("Lista vuota")
    minimo = massimo = valori[0]        # 3 (lettura, 2 assegnamenti)
    for i in range(1, len(valori)):     # n - 1 iter.
        valore = valori[i]              # 2 (lettura, assegnamento)
        if valore < minimo:                    # 1
            minimo = valore             # 1 se vero
        if valore > massimo:            # 1
            massimo = valore            # 1 se vero
    return massimo - minimo             # 1
\end{unipdcode}
\end{minipage}
